\documentclass[10pt]{article}
\usepackage[margin=0.78in]{geometry}
\usepackage{booktabs,array,amsmath,microtype,float}
\usepackage[hidelinks]{hyperref}
\setlength{\parskip}{0.25em}
\setlength{\emergencystretch}{2em}
\title{BidKV-MultiNPU: Ownership-Aware Victim Coordination Under Distributed KV Pressure}
\author{Advisor-authored storyline draft; implementation and experiments remain student-owned}
\date{M0 proposal, August 2026}

\begin{document}
\maketitle
\begin{abstract}
BidKV is an accepted single-device victim-selection policy for KV-pressure LLM
serving.  Running that policy on several accelerators does not by itself create
a multi-device mechanism: a request may have state on several owners, pressure
can be asymmetric, and a locally attractive victim may trigger duplicated
recomputation, collective delay, or cross-device traffic.  This paper draft
poses a separate, falsifiable question: when do ownership and topology make
independent local victim choices globally costly, and can bounded coordination
avoid that cost without global serialization?  No MultiNPU implementation,
deterministic replay, or hardware result is admitted today.  We freeze seven
research questions, a six-link NPU claim chain, local/global/oracle baselines,
conservation oracles, performance and coordination-overhead gates, negative
stops, and a two-week delivery.  Until an ownership counterexample and native
coordination path exist, the extension remains design-level porting/adaptation,
not an NPU mechanism contribution.
\end{abstract}

\section{Introduction}
Paged KV memory allows serving engines to reclaim active request state under
pressure~\cite{Kwon2023}.  BidKV improves the local victim decision by ranking
the capacity a victim frees against its estimated disruption cost
\cite{BidKV2026}.  The accepted manuscript and its NVIDIA results are the
inherited single-device baseline.  They contain no evidence that independent
choices remain correct when one logical request spans or interacts with several
accelerators.

The multi-device extension begins where the local abstraction can fail.  A
decision made from one device's pressure may ignore blocks owned elsewhere,
collective dependencies, asynchronous ownership updates, recomputation that is
charged twice, or transfer induced by moving the decision boundary.  Conversely,
a global coordinator can cost more than it saves.  The contribution therefore
cannot be ``global BidKV.''  It must identify an observable divergence between
local and globally accountable choices, implement the minimum metadata exchange
that repairs it, and predict when coordination should remain off.

The repository currently provides only that research contract.  Historical
single-device result directories remain evidence for inherited BidKV and are
never relabeled MultiNPU.  Dedicated branches identify future carriers but do
not admit code or results into this draft.

\section{Seven-Question Contract}
\begin{table}[H]
\centering\small
\begin{tabular}{p{0.04\linewidth}p{0.89\linewidth}}
\toprule
Q & Frozen answer \\
\midrule
1 & \textbf{Problem.} Under asymmetric per-device KV pressure and shared request ownership, when does independent local victim selection choose a globally expensive victim? \\
2 & \textbf{Importance.} A bad coordinated decision can free the wrong physical capacity, duplicate recomputation, delay collectives, or serialize the serving path. \\
3 & \textbf{Hidden assumption and gap.} Local BidKV assumes that capacity and disruption are attributable at one decision point. Generic multi-resource fairness~\cite{Ghodsi2011} does not supply request/KV ownership or recomputation semantics. \\
4 & \textbf{Mechanism hypothesis.} Exchange a bounded ownership-and-cost receipt only at pressure events, reject stale epochs, and coordinate only when predicted global regret exceeds communication uncertainty; otherwise retain local BidKV. \\
5 & \textbf{Feasibility.} Dedicated parent/runtime branches exist, but exact ownership hooks, replay, and coordination code are pending student work. \\
6 & \textbf{Matched evaluation.} Compare independent local BidKV, the strongest non-BidKV local policy, global largest-first, the treatment, and an offline coordination oracle under identical topology, memory, requests, and restarts. \\
7 & \textbf{Takeaway.} A reusable result is a measured boundary identifying when ownership-aware coordination pays, or a negative result showing that its cost dominates on tested topologies. \\
\bottomrule
\end{tabular}
\caption{Questions are preregistered obligations, not completed claims.}
\end{table}

\section{Ownership Model and Hidden Assumption}
Let $d\in\mathcal D$ index devices.  Request $i$ owns physical KV blocks
$B_{i,d}$, and a local policy observes pressure $P_d$.  Reclaiming $i$ produces
a vector of freed capacity $\mathbf f_i=(f_{i,d})$ and a global disruption
$C_i$ that includes recomputation, synchronization, and any transfer.  A local
score $f_{i,d}/\widehat C_{i,d}$ is safe only if local ownership and local cost
are sufficient statistics for $\mathbf f_i$ and $C_i$.  The project tests that
assumption rather than declaring it false.

A candidate coordination receipt contains request and generation identity,
ownership epoch, per-device blocks, lifecycle state, local pressure, predicted
recompute/transfer cost, and monotonic timestamp.  A stale or incomplete receipt
must not trigger a global action.  The proposed coordinator selects at most the
minimum victim set needed to satisfy explicit per-device deficits and emits an
auditable decision epoch.  This is a design hypothesis, not an implementation.

\section{NPU Six-Link Claim Discipline}
\begin{enumerate}
\item \textbf{Architecture fact.} The selected TP/DP carrier is expected to
distribute request execution and KV ownership, but the precise ownership map and
hook identities are still an M0 deliverable.
\item \textbf{Invalidated assumption.} The candidate invalid assumption is that
each device can price and reclaim victims independently without affecting peers.
It has not been invalidated on the pinned carrier.
\item \textbf{Real counterexample.} None exists.  M0 first requires a
deterministic replay case where the local optimum is globally worse; replay
would still be simulation, not a hardware counterexample.
\item \textbf{Native mechanism.} Ownership receipts and bounded coordination at
named vLLM-HUST/Ascend hook points are pending.
\item \textbf{Predictive boundary.} Coordination should help only when pressure
and ownership are asymmetric, local/global choices diverge, and avoidable global
cost exceeds message, synchronization, and staleness cost.
\item \textbf{Generalization class.} A verified rule may generalize to serving
systems with explicit sharded state ownership; it does not cover replicas with
independent KV state or opaque ownership.
\end{enumerate}
The first four links are incomplete.  BidKV-MultiNPU is therefore explicitly an
\textbf{unverified porting/adaptation study} at this head.

\section{Mechanism Hypothesis}
Coordination is pressure-triggered, not continuously centralized.  Each device
publishes a bounded receipt for its candidate victims.  A decision is accepted
only if all required owners share one epoch and the predicted global regret gap
exceeds uncertainty plus coordination cost.  Otherwise each device uses the
inherited local BidKV policy.  The coordinator neither changes admission nor
compresses KV and never treats a request as released until every owning device
acknowledges its physical block release.

This scope excludes routing, placement, migration policy, cache compression,
load balancing, and new quality surrogates.  It also excludes global fairness as
a primary objective: fairness is a safety metric, while the mechanism claim is
ownership-correct victim coordination.

\section{Evidence Ledger}
\begin{table}[H]
\centering\small
\begin{tabular}{>{\raggedright\arraybackslash}p{0.23\linewidth}>{\raggedright\arraybackslash}p{0.18\linewidth}>{\raggedright\arraybackslash}p{0.50\linewidth}}
\toprule
Asset & Label & Admissible interpretation \\
\midrule
Accepted BidKV paper/results & historical real-online, single device & Establishes the inherited local baseline only. \\
MultiNPU M0 document & design/contract & Freezes ownership questions, replay baselines, invariants, and negative stops. \\
Dedicated feature branches & unreviewed carrier & Identify future integration locations; no mechanism or result is admitted. \\
Deterministic replay & TBD simulation & Must create local/global divergence and pass accounting before runtime work. \\
Multi-NPU campaign & TBD & Required for any latency, throughput, scaling, or overhead statement. \\
\bottomrule
\end{tabular}
\caption{No historical single-device result is promoted to MultiNPU evidence.}
\end{table}

\section{Matched Evaluation Contract}
\subsection{Replay gate}
The deterministic model must include at least two devices, asymmetric pressure,
shared requests, independent local BidKV, the strongest non-BidKV local policy,
global largest-first, and an offline oracle.  It must contain a case where the
local optimum is globally worse.  If the treatment cannot beat local BidKV on
any such case, runtime implementation stops.

\subsection{Runtime identity and correctness}
Any later vLLM-HUST execution uses the project-approved \texttt{.23} official
container.  Host, devices, topology, model, container digest, parent and
submodule heads, graph mode, command, requests, and result root come from a
scheduler/allocation manifest; no fixed device or personal path is allowed.

The oracle requires: one live ownership epoch per decision; no request both
resident and released; exact per-owner freed-block conservation; one attribution
of recompute and transfer cost; no duplicate victim; no stale decision applied;
and token-identical deterministic output versus the inherited local baseline.
Missing receipts fail closed to local BidKV.  Any correctness failure blocks
performance claims.

\subsection{Metrics and gates}
Across at least three independent service starts, report task/request TTFT
median/p95/p99, SLO attainment, throughput, preemptions, recomputed tokens,
cross-device bytes, per-device freed blocks, coordination messages and latency,
decision staleness, local/global divergence, oracle regret, fairness, failures,
and all fallbacks.  Freeze numerical efficacy and overhead thresholds before
opening results.  The minimum qualitative gate is a repeated deployable winner
switch plus oracle headroom over independent local BidKV; an isolated latency
delta is insufficient.

Stop if ownership cannot be observed precisely, coordination cost is comparable
to avoided recomputation, choices rarely diverge, or gains arise from changing
memory, requests, output length, execution mode, or correctness.  Negative
results close this mechanism, not the broader ownership problem.

Week one delivers the ownership map, typed replay, three baselines, oracle, and
invariant tests.  Week two names native hooks and implements a reviewable receipt
path only if replay clears the action-space gate; hardware execution is a later,
separately allocated student deliverable.

\section{Related-Work and Claim Boundary}
The accepted BidKV work supplies the strongest inherited local baseline
\cite{BidKV2026}; its utility signal and results are not new MultiNPU evidence.
PagedAttention documents the serving memory substrate and native reclamation
context~\cite{Kwon2023}.  Dominant Resource Fairness formalizes multi-resource
allocation fairness~\cite{Ghodsi2011}, but it does not bind logical requests to
distributed KV ownership or price duplicated recomputation.  The proposed gap
is therefore not generic global scheduling.  It is the conjunction of typed KV
ownership, pressure-local decisions, and fail-closed coordination economics.
A broader bibliography and direct multi-device victim-selection comparison are
required before submission.

\section{Conclusion}
BidKV-MultiNPU is a separate research question, not a scaling label for an
accepted single-device system.  At the current head there is no replay
counterexample, native mechanism, or hardware evidence.  The project advances
only if ownership-aware coordination produces observable action differences,
passes exact accounting, and beats local BidKV after coordination cost.  A clean
negative boundary is preferable to relabeling inherited results.

\bibliographystyle{plain}
\bibliography{references}
\end{document}
